% Figure — the self-improvement ladder (representative rung: §3.1.5 ASPIRE loop).
\begin{figure}[t]
\centering
\begin{tikzpicture}[
  node distance=6mm and 9mm,
  box/.style={draw=hidden-draw, rounded corners, fill=secAl, align=center,
              font=\footnotesize, minimum height=8mm, text width=15mm, inner sep=3pt},
  core/.style={box, fill=secA},
  io/.style={draw=none, font=\footnotesize\itshape},
  ar/.style={-{Stealth[length=2mm]}, darkgray, line width=0.7pt},
]
  \node[io] (task) {task};
  \node[box, fill=secA, right=of task] (actor) {actor\\ agent};
  \node[box, fill=secB, right=of actor] (eng) {exec.\\ engine (F)};
  \node[box, fill=secC, right=of eng] (mem) {skill\\ memory (M)};
  \node[box, fill=secD, right=of mem] (srch) {evol.\\ search (S)};
  \node[io, right=of srch] (out) {validated skill};

  \draw[ar] (task) -- (actor);
  \draw[ar] (actor) -- (eng);
  \draw[ar] (eng) -- (mem);
  \draw[ar] (mem) -- (srch);
  \draw[ar] (srch) -- (out);
  % feedback loops
  \draw[ar] (eng.south) to[out=-90,in=-90] node[below,font=\scriptsize]{traces} (actor.south);
  \draw[ar] (mem.north) to[out=90,in=90] node[above,font=\scriptsize]{skills $\rightarrow$ future tasks} (actor.north);
\end{tikzpicture}
\caption{The full self-improving loop (\S3.1.5): feedback (F) $+$ memory (M) $+$ search (S)
combine into one open-ended loop---the ASPIRE cell. VLAs (\S3.2) have no such loop.}
\label{fig:architectures}
\end{figure}
